
\section*{7. Quotient by a proper and flat groupoid}
\phantomsection\label{sec:sga3-V-7}
\addcontentsline{toc}{section}{7. Quotient by a proper and flat groupoid}

\medskip
\phantomsection\label{thm:sga3-V-7.1}%
\noindent\textbf{Theorem 7.1.}---%
\footnote{\label{note:sga3-V-34}Editor's note: We mention here the
article of S.~Keel and S.~Mori \hyperref[bib:sga3-V-KM97]{[KM97]}, which establishes the following
theorem. Let \(X\) be an algebraic space of finite type over a locally
Noetherian base \(S\), and let
\(j:R\to X\times_S X\) be a flat groupoid whose stabilizer
\(j^{-1}(\Delta_X)\) is finite over \(X\). There then exists an
algebraic space that is a geometric quotient of \(X\) by \(R\) and a
uniform categorical quotient; moreover, if \(j\) is separated, this
quotient is separated. In particular, if a flat \(S\)-group scheme
\(G\) acts properly on \(X\), with finite stabilizer (that is, if
\[
  G\times_S X\longrightarrow X\times_S X,\qquad
  (g,x)\longmapsto(x,g\cdot x)
\]
is proper and the stabilizer of the diagonal is finite over \(X\)),
then there exists a geometric quotient \(X\to X/G\). For a reductive
\(S\)-group scheme \(G\), this is a result of J.~Koll\'ar \hyperref[bib:sga3-V-Ko97]{[Ko97]}.}%
\emph{Let \(S\) be a locally Noetherian scheme and let}
\SGAThreeDiagram[0.56\textwidth]
  {figures/exp5/p274_theorem71_groupoid.png}
  {The groupoid occurring in \hyperref[thm:sga3-V-7.1]{Theorem~7.1}.}
  {fig:sga3-V-7.1-groupoid}
  {Expos\'e V, component-local p.~24 / combined-reader p.~286 / re-edited header p.~274}
\emph{be a \((\mathbf{Sch}/S)\)-groupoid such that \(d_1\) is proper
and flat, \(X_0\) is quasi-projective over \(S\),}%
\footnote{\label{note:sga3-V-35}Editor's note: This hypothesis on
\(X_0\) is necessary; see \hyperref[note:sga3-V-17]{editor's note~(17)} to
Theorem~\hyperref[thm:sga3-V-4.1]{4.1}.}
\emph{and the morphism
\[
  d:X_1\longrightarrow X_0\times_S X_0
\]
with components \(d_0\) and \(d_1\) is quasi-finite. Then:}
\begin{enumerate}
\item[(i)]\phantomsection\label{thm:sga3-V-7.1-i}%
  \emph{There exists a cokernel \((Y,p)\) of \((d_0,d_1)\) in
  \((\mathbf{Sch}/S)\). Moreover, every such \((Y,p)\) is a cokernel
  of \((d_0,d_1)\) in the category of all ringed spaces.}
\item[(ii)]\phantomsection\label{thm:sga3-V-7.1-ii}%
  \emph{The morphism \(p\) is surjective, open, proper, and of finite
  presentation, and \(Y\to S\) is separated and of finite
  presentation.}%
  \footnote{\label{note:sga3-V-36}Editor's note: TDTE~III,
  \SGARefLink{lem:sga3-V-6.1}{Theorem~6.1}, states that \(Y\to S\) is quasi-projective if \(S\) is
  Noetherian. The editors did not see how to deduce this from the
  local existence of quasi-sections.}
\item[(iii)]\phantomsection\label{thm:sga3-V-7.1-iii}%
  \emph{The morphism
  \[
    X_1\longrightarrow X_0\times_Y X_0
  \]
  with components \(d_0\) and \(d_1\) is surjective.}
\item[(iv)]\phantomsection\label{thm:sga3-V-7.1-iv}%
  \emph{If \((d_0,d_1)\) is an equivalence pair, then
  \(X_1\to X_0\times_Y X_0\) is an isomorphism and \(p\) is
  faithfully flat.}%
  \footnote{\label{note:sga3-V-37}Editor's note: The consequences
  that follow have been made explicit; see \hyperref[bib:sga3-V-Ray67a]{[Ray67a]},
  Theorem~1(iii), and the end of the proof of \hyperref[thm:sga3-V-7.1]{the theorem}. We also
  mention that another proof of \hyperref[thm:sga3-V-7.1]{Theorem~7.1} in the case of an
  equivalence relation, based on the existence of Hilbert schemes,
  is given in \hyperref[bib:sga3-V-AK80]{[AK80]}, Theorem~2.9; see also \hyperref[bib:sga3-V-BLR90]{[BLR90]}, \S8.2,
  Theorem~12. In that case it is furthermore shown that \(Y\to S\)
  is quasi-projective.}
  \emph{Moreover, \((Y,p)\) is a cokernel of \((d_0,d_1)\) in the
  category of sheaves for the fppf topology, and, for every base
  change \(Y'\to Y\), the scheme \(Y'\) is the cokernel of the
  groupoid \(X_\ast\times_Y Y'\) obtained from \(X_\ast\) by the
  base change \(X_0\times_Y Y'\to X_0\).}
\end{enumerate}

\emph{In particular, for every base change \(S'\to S\), the scheme
\[
  Y'=Y\times_S S'
\]
is the cokernel of the \(S'\)-groupoid
\[
  X'_\ast=X_\ast\times_S S'.
\]
Thus, in this case, \textquotedblleft{}formation of the quotient
commutes with base change\textquotedblright.}

Let \((Y,p)\) be the cokernel of \((d_0,d_1)\) in the category of all
ringed spaces. Lemma~\hyperref[sga3:exp5:lemma:1.2]{1.2} shows that,
to prove part~\hyperref[thm:sga3-V-7.1-i]{(i)}, it is enough to show
that every point \(z\in X_0\) has a saturated open neighborhood
\(U_z\) such that, if \(\widetilde d_0\) and
\(\widetilde d_1\) denote the restrictions of \(d_0\) and \(d_1\) to
\[
  d_0^{-1}(U_z)=d_1^{-1}(U_z),
\]
and if \((Q,q)\) denotes the cokernel of
\((\widetilde d_0,\widetilde d_1)\) in the category of ringed spaces,
then \(Q\) is a scheme and \(q\) is a morphism of schemes.

By Lemma~\hyperref[lem:sga3-V-6.1-i]{6.1(i)}, it is therefore enough
to show that every point \(z\in X_0\) has a saturated open
neighborhood \(U_z\) such that the groupoid induced by \(X_\ast\) on
\(U_z\) has a quasi-section. We may even assume that \(z\) is closed
in its fiber over \(S\) (we shall say that \(z\) is \emph{closed
relative to \(S\)}).%
\footnote{\label{note:sga3-V-38}Editor's note: Indeed, if such an open
neighborhood \(U_z\) has been constructed for every point \(z\) that
is closed relative to \(S\), then the union of these \(U_z\) covers
\(X_0\): every fiber over \(S\) of the closed complement is a
Noetherian scheme with no closed points, and is therefore empty.}
The existence of \(U_z\) follows from the lemmas below.
